\documentclass[UTF8,a4paper,11pt]{ctexart}
\usepackage{amsmath,amssymb,graphicx,geometry}
\geometry{margin=23mm}
\setlength{\parindent}{2em}
\setlength{\parskip}{0.35em}
\title{反馈取样位置、时间参数假设与根轨迹的形式化比较}
\author{}
\date{}
\begin{document}
\maketitle
本文分别求解两种反馈取样连线，以及秒时间模型与另一组毫秒归一化参数假设。稳定性均指闭环全部特征根位于开左半平面；阻尼条件逐个检查全部极点。
\section*{一、AP6.5：分别按两种取样连线求解}
设五个增益均为正。第一种连线的局部反馈分别在增益 $3$ 之前、末级 $10/(s-10)$ 之前取样；第二种连线分别在这两个环节之后取样。取样位置改变局部闭环分母，因此分别从节点方程建立模型。

\subsection*{1. 第一种连线：传递函数、稳定范围与阶跃响应}
令 $x$ 为 $1/(s+1)$ 后的信号，$v$ 为 $K_2$ 后、末级动态环节前的信号。由三个节点关系得

\[(s+1)x=K_1(R-K_3x-K_4Y).\]
\[(2s-4)v=K_2(3x-K_5v),\qquad (s-10)Y=10v.\]
记三个组合参数为

\[a=1+K_1K_3>1,\qquad b=K_2K_5-4>-4,\qquad g=30K_1K_2K_4>0.\]
消去 $x,v$，得到题目（a）的闭环传递函数

\[T_a(s)=\frac{Y(s)}{R(s)}=\frac{30K_1K_2}{(s+a)(2s+b)(s-10)+g}.\]
\[D_a(s)=2s^3+As^2+Bs+C.\]
\[A=2a+b-20,\qquad B=ab-20a-10b,\qquad C=g-10ab.\]
对题目（b），三阶劳斯表为

\[\begin{array}{c|cc}s^3&2&B\\s^2&A&C\\s^1&(AB-2C)/A&0\\s^0&C&0\end{array}\]
由劳斯判据得全部正增益下的充要稳定条件

\[A>0,\qquad B>0,\qquad C>0,\qquad AB>2C.\]
下面将其化成便于顺序选参的等价范围。若 $1<a\le10$，则 $A>0$ 给出 $b>20-2a\ge0$；此时 $B=(a-10)b-20a<0$，与 $B>0$ 矛盾。因此稳定时必有 $a>10$，并且

\[B>0\quad\Longleftrightarrow\quad b>\frac{20a}{a-10}>20.\]
这时 $A=2a+b-20>0$ 自动成立。再将 $C>0$ 和 $AB>2C$ 解成关于 $g$ 的不等式，得到完全等价的充要条件

\[a>10,\qquad b>\frac{20a}{a-10}.\]
\[10ab<g<10ab+\frac{(2a+b-20)(ab-20a-10b)}{2}.\]
因此可以先任取 $K_1>0,K_2>0$，再依次选择

\[K_3>\frac9{K_1},\qquad K_5>\frac{4+20a/(a-10)}{K_2}.\]
\[\frac{ab}{3K_1K_2}<K_4<\frac{ab}{3K_1K_2}+\frac{(2a+b-20)(ab-20a-10b)}{60K_1K_2}.\]
上述顺序参数化覆盖所有稳定的正增益组合。下界 $g=10ab$ 对应原点特征根；上界对应

\[D_a(s)=(2s+A)\left(s^2+\frac B2\right),\qquad s=-\frac A2,\ \pm j\sqrt{\frac B2}.\]
为使全部极点位于开左半平面，两个边界均取严格不等号。对题目（c），选取

\[(K_1,K_2,K_3,K_4,K_5)=(1,1,20,300,46).\]
\[a=21,\quad b=42,\quad g=9000,\quad A=64,\quad B=42,\quad C=180.\]
由 $294<K_4<338.8$ 可知所选 $K_4=300$ 位于稳定范围内，且 $AB-2C=2328>0$。相应传递函数与全部极点为

\[T_a(s)=\frac{30}{2s^3+64s^2+42s+180}.\]
\[s=-31.42284533,\qquad -0.28857733\pm j1.66759747\quad(\mathrm{s}^{-1}).\]
对单位阶跃输入，终值定理给出

\[y(\infty)=\lim_{s\to0}T_a(s)=\frac16.\]
若以 $\lambda_i$ 表示这三个简单极点，阶跃响应还可写为

\[y(t)=\frac16+\sum_{i=1}^{3}\frac{30}{\lambda_iD_a^{\prime}(\lambda_i)}e^{\lambda_i t},\qquad t\ge0.\]
共轭极点使响应围绕终值作衰减振荡，衰减率由 $0.28857733\ \mathrm{s}^{-1}$ 决定。数值状态空间计算与留数表达式得到相同的阶跃曲线。

\begin{figure}[htbp]\centering
\includegraphics[width=\linewidth]{\detokenize{AP6_5_step.png}}
\caption{第一种取样连线在所选正增益下的单位阶跃响应，终值为 1/6。}
\end{figure}
\subsection*{2. 第二种连线：传递函数与正增益稳定集}
令 $z$ 为增益 $3$ 后的信号。两条局部反馈分别取自 $z$ 和输出 $Y$，由框图得

\[(s+1)z=3K_1(R-K_3z-K_4Y).\]
\[(2s-4)(s-10)Y=10K_2(z-K_5Y).\]
记

\[c=1+3K_1K_3>1,\qquad d=40+10K_2K_5>40,\qquad g=30K_1K_2K_4>0.\]
则题目（a）的闭环传递函数为

\[T_b(s)=\frac{30K_1K_2}{(s+c)(2s^2-24s+d)+g}.\]
\[D_b(s)=2s^3+(2c-24)s^2+(d-24c)s+cd+g.\]
对题目（b），若该三阶系统稳定，系数正性首先要求

\[c>12,\qquad d>24c.\]
劳斯首列还要求下式为正：

\[\Delta=(2c-24)(d-24c)-2(cd+g).\]
展开并利用 $c>12,d>40,g>0$，得

\[\Delta=-24d-48c(c-12)-2g<0.\]
由此得到矛盾，所以题目规定的正增益稳定集为空。对题目（c），稳定范围内的增益组合及相应稳定阶跃响应也构成空集。其原因可以直接从系数关系看出：增大末级输出反馈会提高 $d$，同时使常数项含有 $cd$，两者对劳斯行列式合成的贡献为 $-24d$；外环正增益又贡献 $-2g$。

两种连线的比较结论是：第一种连线具有上面给出的非空稳定参数域；第二种连线在五个增益全正的条件下无稳定参数。每一结论均对应其实际反馈取样位置。

\section*{二、P7.11：秒时间模型的完整解答}
\subsection*{1. 由多环反馈建立模型}
令 $s$ 为以秒为基准的拉普拉斯变量，电机系数按题给数值 $K_T/(LJ)=2$ 解释。采用 $R/L=2000\ \mathrm{s}^{-1}$、$K_b=0.4$、$r=0.2$、$0.5K_1=1$、$K_p=1$。设 $\Omega$ 为轴角速度、$V$ 为放大器输出，反电动势局部反馈给出

\[(s^2+2000s+0.8)\Omega=2V,\qquad M_s(s)=\frac{\Omega}{V}=\frac2{s^2+2000s+0.8}.\]
转速反馈与位置反馈满足

\[(0.001s+1)V=K_a\left(\frac{R-Y}{0.001s+1}-K_2\Omega\right),\qquad sY=0.2\Omega.\]
消去 $V,\Omega$，得到位置闭环传递函数

\[T_s(s)=\frac{0.4K_a}{Q_s(s)}.\]
\[Q_s(s)=s(0.001s+1)^2(s^2+2000s+0.8)+2K_aK_2s(0.001s+1)+0.4K_a.\]
\subsection*{2. 问（a）：固定转速反馈增益的根轨迹}
取 $K_2=10$，将特征方程写成 $D_s(s)+K_aN_s(s)=0$，其中

\[D_s(s)=s(0.001s+1)^2(s^2+2000s+0.8).\]
\[N_s(s)=0.02s^2+20s+0.4.\]
根轨迹的五个开环极点为

\[s=0,\quad -1000\ (\text{二重}),\quad -1000\pm\sqrt{999999.2}.\]
\[s\approx0,\ -1000\ (\text{二重}),\ -0.00040000008,\ -1999.99959999992.\]
两个有限零点为

\[s=-500\pm\sqrt{249980}\approx-0.020000400016,\ -999.979999599984.\]
三条无穷远渐近线的交点和方向为

\[\sigma_a=\frac{-4000-(-1000)}{5-2}=-1000,\qquad \theta_a=60^\circ,180^\circ,300^\circ.\]
实轴根轨迹的三个区段为

\[(-\infty,-1999.99959999992),\quad(-999.979999599984,-0.020000400016).\]
\[(-0.00040000008,0).\]
箭头对应 $K_a$ 由零增大。在二重极点 $-1000$ 附近，令 $s=-1000+\delta$，首项满足 $\delta^2\sim-0.4K_a/999.9992$，故两支沿 $\pm90^\circ$ 出发。原点和慢实极点相向运动并在实轴分离；最左极点沿负实轴趋于负无穷。

由 $K_a=-D_s/N_s$ 及 $D_s^{\prime}N_s-D_sN_s^{\prime}=0$，筛选正增益得到四个实轴重根位置：

\[(s,K_a)\approx(-0.000201010117,0.000202025294).\]
\[(s,K_a)\approx(-999.959999200,4.000156805).\]
\[(s,K_a)\approx(-0.0398009661,7.919641968).\]
\[(s,K_a)\approx(-422.644187354,19246.5633593).\]
两对共轭分支分别在约 $K_a=4.00016$ 和 $7.91964$ 处进入实轴；随后其中两条实轴分支在 $K_a\approx19246.56$ 处再离开实轴。最终两支到达有限零点，三支沿上述渐近方向趋于无穷远。虚轴交点为

\[K_a\approx299999.0799679,\qquad s=\pm j1414.20677404\ \mathrm{s}^{-1}.\]
该根轨迹的正增益稳定范围为

\[0<K_a<299999.0799679.\]
\subsection*{3. 问（b）：全部极点的阻尼条件与选参}
对非零极点定义 $\zeta=-\operatorname{Re}s/|s|$，负实极点取 $\zeta=1$。边界 $\zeta=0.6$ 在上半平面表示为

\[s=-\sigma+j\frac43\sigma,\qquad \sigma>0.\]
将边界极点代入增益关系并令增益为正实数：

\[\operatorname{Im}\!\left[-\frac{D_s(-\sigma+j4\sigma/3)}{N_s(-\sigma+j4\sigma/3)}\right]=0,\qquad K_a>0.\]
求解得到三个正增益边界及对应的上半平面极点：

\[K_a\approx0.0005715445313,\qquad s\approx-0.000202857622+j0.000270476829.\]
\[K_a\approx2.799590930113,\qquad s\approx-0.014197853787+j0.018930471717.\]
\[K_a\approx40981.2562875,\qquad s\approx-359.683773567+j479.578364756.\]
由全部边界划分增益轴，并逐区间检查五个极点，得到题目所求的完整范围

\[0<K_a\le0.0005715445313\quad\text{或}\quad2.799590930113\le K_a\le40981.2562875.\]
小增益段的慢共轭极点首先离开阻尼扇区；继续提高增益后又进入扇区，随后快速分支在第三个边界离开。这一过程使可行集分成两段。选取第二段内的 $K_a=1000$，在 $K_2=10$ 下有

\[s\approx-2009.85357678,\ -999.97997956,\ -980.01241949.\]
\[s\approx-10.13398442,\ -0.02003975721\quad(\mathrm{s}^{-1}).\]
五个极点均为负实数，因而全部阻尼比为 $1$。最慢极点的时间常数约为 $49.90\ \mathrm{s}$。阻尼条件控制极点角度，而响应时间还与极点实部有关。

\subsection*{4. 问（c）：固定放大器增益时的根轨迹}
固定所选 $K_a=1000$，对 $K_2$ 写成

\[D_{2s}(s)+K_2N_{2s}(s)=0,\quad D_{2s}(s)=D_s(s)+400,\quad N_{2s}(s)=2s^2+2000s.\]
开环极点为

\[s\approx-1999.99970000,\qquad-1000.00020000\pm j0.63245550.\]
\[s\approx0.000049999855\pm j0.44721344736.\]
有限零点为 $0,-1000$；三条无穷远渐近线仍交于 $-1000$，方向仍为 $60^\circ,180^\circ,300^\circ$。实轴轨迹区段为

\[(-\infty,-1999.99970000),\qquad(-1000,0).\]
上半平面的快速极点出射角约为 $0.01812^\circ$，慢极点出射角约为 $179.95516^\circ$，下半平面关于实轴对称。慢共轭分支随 $K_2$ 增大先向左穿过虚轴；实轴重根位置为

\[(s,K_2)\approx(-999.367543962,0.632655279).\]
\[(s,K_2)\approx(-0.447513898,0.893927303).\]
\[(s,K_2)\approx(-422.649661750,192.450678405).\]
两次虚轴交叉给出

\[K_2\approx0.0000999998100001,\qquad s=\pm j0.447213483697\ \mathrm{s}^{-1}.\]
\[K_2\approx3000.000766667,\qquad s=\pm j1414.213821646\ \mathrm{s}^{-1}.\]
因此完整非负增益稳定范围为

\[0.0000999998100001<K_2<3000.000766667.\]
若同时保持问（b）的全部极点阻尼要求，则由同一阻尼边界求得

\[0.536540544856\le K_2\le409.800363717.\]
\[s\approx-0.268436179553\pm j0.357914906071\quad\text{（下边界）}.\]
\[s\approx-359.687584229\pm j479.583445638\quad\text{（上边界）}.\]
所选 $K_2=10$ 同时位于稳定范围与阻尼可行范围内。

\section*{三、P7.11：独立的毫秒归一化参数假设}
\subsection*{1. 模型约定与问（a）的根轨迹}
现在另行规定一组以毫秒为基准的归一化参数：两个一阶环节均为 $1/(p+1)$，电机局部闭环为 $2/(p^2+2p+0.8)$，其余位置和转速反馈按此归一化模型取值。这组规定保留了电机系数的数值 $2$，对应另一组动态参数。其闭环模型为

\[T_m(p)=\frac{0.4K_a}{Q_m(p)}.\]
\[Q_m(p)=p(p+1)^2(p^2+2p+0.8)+2K_aK_2p(p+1)+0.4K_a.\]
问（a）固定 $K_2=10$，有

\[D_m(p)=p(p+1)^2(p^2+2p+0.8),\qquad N_m(p)=20p^2+20p+0.4.\]
五个开环极点与两个有限零点分别为

\[p=0,\quad-1\ (\text{二重}),\quad-1\pm\sqrt{0.2}.\]
\[p\approx0,\ -1\ (\text{二重}),\ -0.5527864045,\ -1.4472135955.\]
\[p=\frac{-1\pm\sqrt{0.92}}2\approx-0.02041684767,\ -0.97958315233.\]
三条渐近线交于 $(-4+1)/3=-1$，方向为 $60^\circ,180^\circ,300^\circ$。实轴轨迹区段为

\[(-\infty,-1.4472135955),\quad(-0.97958315233,-0.5527864045),\quad(-0.02041684767,0).\]
二重极点 $-1$ 沿 $\pm90^\circ$ 出发；原点分支沿实轴趋向零点 $-0.02041685$，最左极点沿实轴趋于负无穷。正增益实轴重根为

\[(p,K_a)\approx(-0.95883095755,0.000827425545).\]
\[(p,K_a)\approx(-0.74889111232,0.001924022164).\]
二重极点出发的共轭分支先进入实轴，其中一支到达 $-0.97958315$，另一支与由 $-0.55278640$ 出发的分支相遇后再成为共轭分支。虚轴交点及正增益稳定范围为

\[K_a\approx0.379845535878,\qquad p=\pm j1.66900916656\ \mathrm{ms}^{-1}.\]
\[0<K_a<0.379845535878.\]
\subsection*{2. 问（b）：阻尼范围与选参}
代入阻尼边界 $p=-\sigma+j4\sigma/3$，求解 $\operatorname{Im}[-D_m(p)/N_m(p)]=0$ 并保留正增益，得到唯一正增益阻尼边界

\[K_a\approx0.0350433867877,\qquad p\approx-0.522465858766+j0.696621145021.\]
全部极点满足 $\zeta\ge0.6$ 的正增益范围为

\[0<K_a\le0.0350433867877.\]
取 $K_a=0.02$、$K_2=10$，全部极点为

\[p\approx-1.82972268546,\ -0.97936620609,\ -0.00682019053.\]
\[p\approx-0.59204545896\pm j0.55141882324\quad(\mathrm{ms}^{-1}).\]
最小阻尼比约为 $0.73176816$，满足要求。最慢极点的时间常数约为 $146.62\ \mathrm{ms}$。两组假设的阻尼可行集具有不同拓扑：秒模型分成两段，而本组假设只有从零开始的一段正增益区间。

\begin{figure}[htbp]\centering
\includegraphics[width=\linewidth]{\detokenize{P7_11_Ka_comparison.png}}
\caption{固定 K₂=10 的两组根轨迹及全部极点的最小阻尼比。叉号为极点，圆圈为零点，斜虚线为阻尼边界，箭头表示 Kₐ 增大；两组各自使用标明的时间单位。}
\end{figure}
\subsection*{3. 问（c）：固定所选放大器增益的根轨迹}
固定 $K_a=0.02$，特征方程为

\[[D_m(p)+0.008]+K_2[0.04p(p+1)]=0.\]
开环极点为

\[p\approx-1.47351276974,\ -0.48765621667,\ -0.01048655498.\]
\[p\approx-1.01417222930\pm j0.18200145256\quad(\mathrm{ms}^{-1}).\]
有限零点为 $0,-1$；渐近线交点为 $-1$，方向为 $60^\circ,180^\circ,300^\circ$。实轴轨迹区段为

\[(-\infty,-1.47351276974),\quad(-1,-0.48765621667),\quad(-0.01048655498,0).\]
上半平面复极点的出射角约为 $2.01078^\circ$，下半平面对称。最左实分支趋于负无穷，另外两条实分支分别趋于零点 $-1$ 和 $0$，共轭分支向右上、右下延伸并趋于 $\pm60^\circ$ 渐近线。非负增益范围没有实轴分离或会合点。

非负增益根轨迹仅有一个虚轴交叉：

\[K_2\approx189.996224700,\qquad p=\pm j1.67309526607\ \mathrm{ms}^{-1}.\]
因此稳定范围与同时满足全部阻尼要求的范围分别为

\[0\le K_2<189.996224700.\]
\[0\le K_2\le17.3129092328.\]
阻尼上界的共轭极点为

\[p\approx-0.522359429564\pm j0.696479239419\quad(\mathrm{ms}^{-1}).\]
所选 $K_2=10$ 位于两范围内。与秒模型相比，这组参数在 $K_2=0$ 时已经稳定，故稳定范围从零开始；秒模型则需先通过一个正的转速反馈阈值。

\begin{figure}[htbp]\centering
\includegraphics[width=\linewidth]{\detokenize{P7_11_K2_comparison.png}}
\caption{固定各自所选 Kₐ 后随 K₂ 变化的根轨迹及全部极点的最小阻尼比。叉号为极点，圆圈为零点，斜虚线为阻尼边界，箭头表示 K₂ 增大。}
\end{figure}
\section*{四、边界方程、完整性与正确单位变换}
\subsection*{1. 统一形式与全部边界的求法}
为给出前述数值范围的可核查方程，将两种模型统一写成

\[Q(x)=D_e(x)+K_a[uK_2x(x+1)+\beta].\]
\[D_e(x)=x(x+1)^2(x^2+2x+e).\]
对秒模型取 $x=0.001s$，有 $e=8\times10^{-7},u=2\times10^{-6},\beta=4\times10^{-10}$；对独立毫秒参数假设取 $x=p$，有 $e=0.8,u=2,\beta=0.4$。该正比例缩放保留极点阻尼比。

固定 $K_2=10$，记 $r_0=\beta/(10u)$，在阻尼线上取 $x=-t+j4t/3$。由增益必须为实数得

\[\operatorname{Im}\!\left[-\frac{D_e(x)}{10u[x(x+1)+r_0]}\right]=0,\qquad t>0.\]
除去非零常数及原点因子，得到六次方程

\[F(t)=6875t^6-38700t^5+A_4t^4+A_3t^3+A_2t^2-2916r_0(1+e)t+729er_0=0.\]
\[A_4=38025+5625e-7011r_0,\qquad A_3=-20250-4050e+9072r_0.\]
\[A_2=4050+2025e+4455r_0+891er_0.\]
对全部正实根回代

\[K_a=-\frac{D_e(x)}{10u[x(x+1)+r_0]},\qquad x=-t+j4t/3,\]
并保留正实增益，就分别得到秒模型的三个边界和独立毫秒假设的一个边界。相邻边界之间极点连续变化；在每个区间选一个增益检查全部五根，即确定整段的阻尼条件。由于 $K_a>0$ 时常数项为正，原点不属于其特征根，因而没有遗漏穿过原点的边界。

固定 $K_a$ 时，记 $h_0=\beta K_a$。同样方法将 $D_e+h_0$ 与 $x(x+1)$ 配对，得到阻尼边界方程

\[G(t)=6875t^6-38700t^5+(38025+5625e)t^4-(20250+4050e)t^3+(4050+2025e)t^2+1458h_0t-729h_0=0.\]
\[K_2=-\frac{D_e(x)+h_0}{uK_ax(x+1)},\qquad x=-t+j4t/3.\]
该方程的正实根回代后，秒模型给出两个正增益阻尼边界，独立毫秒假设给出一个，形成前述完整的 $K_2$ 阻尼范围。

稳定边界也可直接求得。令 $x=jw$，$X=w^2>0$，由特征方程的实部和虚部分别为零，消去 $K_2$ 得

\[X^3-(1+e)X^2-(2+e)X+\beta K_a=0.\]
\[K_2=\frac{-X^2+(5+e)X-e}{uK_a},\qquad w=\sqrt X.\]
秒模型代入 $K_a=1000$ 得到两个正 $K_2$ 交点；独立毫秒假设代入 $K_a=0.02$，其中一个交点的增益约为 $-19.586674576$，另一个为 $189.996224700$，故非负增益中只有后者。结合各区间的极点所在半平面即可得到所列稳定范围。

固定 $K_2=10$ 时，用同样消元方法可写成

\[X^3-(1+e+r_0)X^2+[-(2+e)+r_0(5+e)]X-r_0e=0.\]
\[K_a=\frac{-X^2+(5+e)X-e}{10u}.\]
两组参数均仅保留一个正增益虚轴交点，分别给出前述 $K_a$ 稳定上界。边界增益处存在纯虚特征根，所以稳定区间的对应端点取开；阻尼边界本身满足 $\zeta=0.6$，所以取闭。$K_a=0$ 时原点极点仍存在，故两个阻尼可行集均写成严格正增益区间。

\subsection*{2. 同一秒模型换用毫秒标度时发生什么}
对同一个物理模型，设 $t_{\mathrm{ms}}=1000t_{\mathrm{s}}$。拉普拉斯变量满足 $p=0.001s$，即 $s=1000p$。将这一关系代入秒模型并除以非零常数 $10^9$，得到严格恒等式

\[\frac{Q_s(1000p)}{10^9}=p(p+1)^2(p^2+2p+8\times10^{-7})+2\times10^{-6}K_aK_2p(p+1)+4\times10^{-10}K_a.\]
相应位置闭环传递函数为

\[T_s(1000p)=\frac{4\times10^{-10}K_a}{Q_s(1000p)/10^9}.\]
因此，真正的单位变换将秒模型中的全部极点数值除以 $1000$，并同时改变多项式内与时间量纲有关的系数。若 $s_i$ 为秒模型的极点，则

\[p_i=\frac{s_i}{1000},\qquad -\frac{\operatorname{Re}p_i}{|p_i|}=-\frac{\operatorname{Re}s_i}{|s_i|}.\]
根的左右半平面位置、阻尼比、以同一物理增益定义的 $K_a,K_2$ 可行范围以及实际输出随物理时间的变化全部保持一致；只是时间数值与极点数值按相反比例改变。以 $K_a=1000,K_2=10$ 为例，真正换元后的五根为

\[p\approx-2.00985357678,\ -0.99997997956,\ -0.98001241949.\]
\[p\approx-0.01013398442,\ -0.00002003975721\quad(\mathrm{ms}^{-1}).\]
其最慢时间常数为 $49.90\ \mathrm{s}=49900\ \mathrm{ms}$；这与同一模型的秒表示完全一致。相比之下，独立毫秒假设的电机常数项为 $0.8$、速度反馈系数为 $2$、位置反馈系数为 $0.4$，而真正换元后分别为 $8\times10^{-7}$、$2\times10^{-6}$、$4\times10^{-10}$。这些系数的差别改变了系统本身，因此两组假设得到不同的根轨迹和增益范围。

比较时应分别理解三个对象：明确采用秒变量的原数值模型；另行规定并保留数值 $2$ 的毫秒归一化参数模型；以及由前者作 $s=1000p$ 得到的同一系统的毫秒表示。前两个是不同参数模型，第一与第三是同一模型的两种单位表示。

\end{document}